\documentclass[11pt]{article}
\usepackage[a4paper,margin=1in]{geometry}
\usepackage{fontspec}
\usepackage[boldfont=false]{xeCJK}
\setCJKmainfont{FandolSong-Regular.otf}
\usepackage{amsmath}
\usepackage{amssymb}
\usepackage{array}
\usepackage{booktabs}
\usepackage{graphicx}
\usepackage{adjustbox}
\usepackage{enumitem}
\setlist[itemize]{topsep=2pt,partopsep=0pt,parsep=0pt,itemsep=2pt,leftmargin=1.4em}
\setlist[enumerate]{topsep=2pt,partopsep=0pt,parsep=0pt,itemsep=2pt,leftmargin=1.6em}
\usepackage{parskip}
\usepackage{fvextra}
\fvset{breaklines=true,breakanywhere=true,fontsize=\small}
\setlength{\parindent}{0pt}

\begin{document}

\begin{center}
  {\LARGE\bfseries Plant nutrition}\\[0.35em]
  {\large Igcse Biology}
\end{center}
\vspace{0.6em}

\section*{Photosynthesis}

\textbf{Photosynthesis} 光合作用 is the process by which plants make \textbf{carbohydrates} 碳水化合物 from simple \textbf{raw materials} 原料, using \textbf{energy} 能量 from light. It is how plants feed themselves.

\subsection*{The word equation}

\[
\text{carbon dioxide} + \text{water} \rightarrow \text{glucose} + \text{oxygen}
\]

In words: plants take in \textbf{carbon dioxide} 二氧化碳 and water, and make \textbf{glucose} 葡萄糖 and \textbf{oxygen} 氧气. This reaction can only happen in the light, and only where there is \textbf{chlorophyll} 叶绿素.

\textbf{(Supplement)} The balanced chemical equation is:

\[
6\text{CO}_2 + 6\text{H}_2\text{O} \rightarrow \text{C}_6\text{H}_{12}\text{O}_6 + 6\text{O}_2
\]

\subsection*{What chlorophyll does}

Chlorophyll is a green \textbf{pigment} 色素 found in the \textbf{chloroplasts} 叶绿体. It traps light and transfers the energy from light into energy stored in chemicals. That chemical energy is then used to build carbohydrates.

\subsection*{What the plant does with the glucose}

\begin{center}
\adjustbox{max width=\linewidth}{%
\begin{tabular}{ll}
\toprule
\textbf{Use} & \textbf{What happens} \\
\midrule
store & changed to \textbf{starch} 淀粉, an energy store that does not dissolve \\
build & made into \textbf{cellulose} 纤维素 to build \textbf{cell walls} 细胞壁 \\
release energy & broken down in \textbf{respiration} 呼吸作用 to release energy \\
transport & changed to \textbf{sucrose} 蔗糖 and moved around the plant in the \textbf{phloem} 韧皮部 \\
attract animals & made into \textbf{nectar} 花蜜 to attract \textbf{insects} 昆虫 for \textbf{pollination} 传粉 \\
\bottomrule
\end{tabular}}
\end{center}

\subsection*{Minerals from the soil}

Plants also take up mineral \textbf{ions} 离子 from the soil through their roots:

\begin{itemize}
\item \textbf{nitrate ions} 硝酸根离子 — needed to make \textbf{amino acids} 氨基酸 (and so proteins).

\item \textbf{magnesium ions} 镁离子 — needed to make chlorophyll.

\end{itemize}

A plant short of nitrate grows poorly with weak stems; a plant short of magnesium has yellow leaves (it cannot make enough chlorophyll).

\subsection*{The rate of photosynthesis}

The \textbf{rate} 速率 of photosynthesis depends on three things:

\begin{itemize}
\item \textbf{light intensity} 光照强度 — more light gives a faster rate, up to a point.

\item carbon dioxide \textbf{concentration} 浓度 — more carbon dioxide gives a faster rate.

\item \textbf{temperature} 温度 — a warmer temperature is faster, until it gets so hot that the enzymes are denatured.

\end{itemize}

\textbf{(Supplement)} At any moment, the one factor that is in shortest supply holds back the rate. This is called the \textbf{limiting factor} 限制因素. For example, on a dull day light intensity is usually the limiting factor; on a bright day carbon dioxide may be.

\subsection*{Investigating photosynthesis}

To show that a plant needs light, chlorophyll and carbon dioxide, you test a leaf for starch after taking one of them away. A \textbf{control} 对照 keeps every other condition the same, so the test is fair. (First leave the plant in the dark to use up its starch.) A leaf with green and white parts shows starch only in the green parts, which have chlorophyll.

You can also watch \textbf{gas exchange} 气体交换 in a water plant using hydrogencarbonate \textbf{indicator} 指示剂, which changes colour as the carbon dioxide level changes. In the light the plant takes in carbon dioxide for photosynthesis; in the dark it gives out carbon dioxide from respiration.

\section*{Leaf structure}

A leaf is well \textbf{adapted} 适应 for photosynthesis. Most leaves are broad, with a large \textbf{surface area} 表面积 to catch plenty of light, and thin, so that gases and light quickly reach all the cells.

The middle of the leaf, the \textbf{mesophyll} 叶肉, is where most photosynthesis happens. These are the main parts of a leaf:

\begin{center}
\adjustbox{max width=\linewidth}{%
\begin{tabular}{ll}
\toprule
\textbf{Structure} & \textbf{Job / how it helps photosynthesis} \\
\midrule
\textbf{cuticle} 角质层 & a clear waxy layer on top; it reduces water loss but still lets light through \\
upper \textbf{epidermis} 表皮 & a clear layer of cells with no chloroplasts; it lets light pass to the cells below \\
\textbf{palisade mesophyll} 栅栏叶肉 & tall, column-shaped cells packed with chloroplasts, near the top; they do most of the photosynthesis \\
\textbf{spongy mesophyll} 海绵叶肉 & rounded cells with \textbf{air spaces} 气腔 between them, so gases can move easily \\
\textbf{stomata} 气孔 & tiny holes, mostly on the lower surface, that let carbon dioxide in and oxygen out \\
\textbf{guard cells} 保卫细胞 & a pair of cells around each stoma; they open and close it \\
lower epidermis & a thin layer that holds the stomata \\
\textbf{vascular bundles} 维管束 & contain \textbf{xylem} 木质部 (brings water to the leaf) and phloem (carries sugars away) \\
\bottomrule
\end{tabular}}
\end{center}

\section*{Exam tips}

\begin{itemize}
\item Learn the word equation exactly: carbon dioxide + water → glucose + oxygen (in light, with chlorophyll).

\item Chlorophyll is in the chloroplasts; it does not get used up — it transfers light energy.

\item Link each factor to the rate: more light, more carbon dioxide, warmer (but not too hot) → faster, until a \textbf{limiting factor} stops further increase.

\item For leaf structure, always link the part to its job: broad and thin to catch light; palisade cells full of chloroplasts; stomata and air spaces for gas exchange.

\item Nitrate ions → amino acids; magnesium ions → chlorophyll.

\end{itemize}


\end{document}
